Here, we provide additional details on our experimental setup.
%Our codebase itself has has been included as part of the supplementary material.

\subsection{Model architecture}
\label{app:architecture}

We use architectures from the GPT-2 family~\citep{radford2018improving} as implemented by HuggingFace~\citep{wolf-etal-2020-transformers}
~\footnote{\url{https://huggingface.co/docs/transformers/model_doc/gpt2}} .
Specifically, we consider the following set of configurations.
\begin{center}
    \begin{tabular}{lrrrr}
    \toprule
        Model & Embedding size & \#Layers & \#Heads & (Total parameters) \\
    \midrule
        Tiny & 64 & 3 & 2 & 0.2M \\
        Small & 128 & 6 & 4 & 1.2M \\
        \textbf{Standard} & 256 & 12 & 8 & 9.5M \\
    \bottomrule
    \end{tabular}
\end{center}
We use the Standard model for the bulk of our experiments and only consider the smaller models for the capacity explorations in Section~\ref{sec:factors} and Appendix~\ref{app:capacity}.
Since we train on each input once (we sample new inputs at each training step), overfitting to the training data is not an issue. Therefore, we set the Dropout probability to 0.

Out of the box, these models take as input a sequence of vectors in embedding space and output a sequence of vectors in the same space.
However, the tasks we study are functions from a lower dimensional vector space (e.g., 10-50 dimensions) to a scalar value.
Thus, in order to use a prompt such as $x_1, f(x_1), x_2, f(x_2)\ldots$, we need to map $x_i$s and $f(x_i)$s to vectors in embedding space.
We do so by first turning the scalars $f(x_i)$ into vectors of the same dimension as $x_i$ by appending 0s and then applying a learnable linear transformation to map all these vectors into the embedding space.
Finally, we map the model output into a scalar value through a dot product with a learnable vector.

We treat the prediction of the model at the position corresponding to $x_i$ (that is absolute position $2i-1$) as the prediction of $f(x_i)$.
Due to the structure of these models, this prediction only depends on $(x_j, f(x_j))$ for $j<i$ and $x_i$.
We ignore the model predictions at positions corresponding to $f(x_i)$.

\subsection{Training}
\label{app:training}
Each training prompt is produced by sampling a random function $f$ from the function class we are training on, then sampling inputs $x_i$ from the isotropic Gaussian distribution $N(0, I_d)$ and constructing a prompt as $(x_1, f(x_1),\ldots,x_k, f(x_{k}))$.
Given a prompt, we obtain model predictions $\hat y_i$ (meant to approximate $f(x_i)$) for each input, and compute the loss
\[
    \frac{1}{k} \sum_{i=1}^k \left(\hat{y_i}-f(x_i)\right)^2.
\]
At each training step, we average the loss over a batch of randomly generated prompts (with different functions and prompt inputs), and perform an update step.
We use the Adam optimizer \citep{kingma2014adam}, and train for 500,000 total steps with a batch size of 64. We use a learning rate of $10^{-4}$ for all function classes and models. 


\paragraph{Curriculum learning.} 
To accelerate training, we start by training on prompt inputs $x_i$ lying in a smaller dimensional subspace, and with fewer inputs per prompt, and gradually increase the subspace dimension and number of prompt inputs.
Specifically, we zero out all but the first $d_\text{cur}$ coordinates of $x_i$, sample prompts of size $k_\text{cur}$ and leave the rest of the training process the same.
We use the same schedule for all training runs for the function classes of linear functions and sparse linear functions, starting with $d_\text{cur} = 5,~ k_\text{cur} = 11$, and increasing $d_\text{cur}$ and $k_\text{cur}$ by $1$ and $2$ respectively, every 2000 steps, until $d_\text{cur} = d$, $k_\text{cur} = 2d+1$. We use a slightly different schedule for 2 layer neural networks and decision trees as we want prompts with more inputs for these function classes.
For these classes, we start with $d_\text{cur} = 5,~ k_\text{cur} = 26$, and increase $d_\text{cur}$ and $k_\text{cur}$ by $1$ and $5$ respectively, every 2000 steps, until $d_\text{cur} = d$, $k_\text{cur} = 5d+1$.

Overall, with curriculum, a training prompt $(x_1, f(x_1), \ldots, x_{k_\text{cur}}, f(x_{k_\text{cur}})$ is generated by sampling a random function $f$ from the function class,  drawing inputs $x_i$ by sampling i.i.d.  from $N(0, I_d)$ and zeroing out all but the first $d_\text{cur}$ coordinates. Given model predictions $\hat{y_i}$, the loss is given by 
\[
    \frac{1}{k_\text{cur}} \sum_{i=1}^{k_\text{cur}} \left(\hat{y_i}-f(x_i)\right)^2.
\]

\paragraph{Sampling random functions.} For the class of linear functions, we sample random function $f(x) = w^\top x$ by drawing $w \sim N(0, I_d)$. For our main setting (Section \ref{sec:linear_functions} and \ref{sec:distribution_shifts}), we set $d = 20$.

For the class of two-layer neural networks, we sample $f(x) = \sum_{i = 1}^r \alpha_i \sigma(w_i^\top x)$, where $\alpha_i$s and $w_i$s are drawn i.i.d. from $N(0, 2/r)$ and $N(0, I_d)$ respectively.  We set $d = 20$ and $r = 100$.

For the class of $k$-sparse linear functions, we sample $f(x) = w^\top x$ by drawing $w \sim N(0, I_d)$ and zeroing out all but $k$ coordinates of $w$ chosen uniformly at random from the first $d_\text{cur}$ coordinates (as defined in the curriculum learning description above).
We set $d = 20$ and $k = 3$. 

For the class of decision trees, the random function $f$ is represented by a decision tree of depth $4$ (with 16 leaf nodes), with 20 dimensional inputs. Each non-leaf node of the tree is associated with a coordinate selected uniformly at random from $\{1, 2, \ldots, d\}$, and each leaf node is associated with a value drawn randomly from $N(0, 1)$. To evaluate $f$ on an input $x$, we traverse the tree starting from the root node, and go to the right child if the coordinate associated with the current node is positive and go to the left child otherwise. $f(x)$ is given by the value associated with the leaf node reached at the end.

\paragraph{Computational resources.}
We train using a single NVIDIA GeForce RTX 3090 GPU and most training runs  take 5-20 hours depending on model size and context length. For instance, for the class of linear functions, training the standard model takes 17 hours for $d = 50$, 7 hours for $d = 20$ and 5.5 hours for $d = 10$. For decision trees, training the standard model takes 17 hours. The time it takes for decision trees and $50$ dimensional linear functions is higher due to larger context lengths (we train for $d$ dimensional linear functions with $2d+1$ input-output pairs per prompt).

\subsection{Baselines}
\label{app:baselines}
\paragraph{Least squares.} Minimum norm least squares is the optimal estimator for the linear regression problem. Given a prompt $P = (x_1, y_1, \ldots, x_k, y_k, x_{\text{query}})$, let $X$ be a $k \times d$ matrix with row $i$ given by $x_i$, and let $y$ be a $k$ dimensional vector with the $i^{\text{th}}$ entry $y_i$. Set $\hat{w}^T  = X^{+} y$, where $X^{+}$ denotes the Moore-Penrose pseudoinverse of $X$. The estimator predicts $M(P) = \hat{w}^T x_{\text{query}}$.

\paragraph{Averaging estimator.} This corresponds to $M(P) = \hat{w}^T x_{\text{query}}$ where $\hat{w} = \frac{1}{k} \sum_{i=1}^k x_i y_i$. This estimator is consistent (yet sub-optimal) when $x_i$s are drawn from $N(0, I_d)$. Unlike least squares, this estimator does not involve an inverse computation, and might be easier for a model to encode.

\paragraph{Nearest neighbors.} This corresponds to setting $M(P) = \frac{1}{n} \sum_{i \in S} y_i$. Here, $S$ is the set of indices of the $n$ nearest neighbors of $x_{\text{query}}$ among $x_1$ to $x_k$. For $k < n$, we average over all the $y_i$s from $1$ to $k$, and for $k=0$, we set $M(P) = 0$. 
We consider the nearest neighbors baselines as it might be easier for a Transformer model to encode using self-attention compared to least squares.


\paragraph{Lasso.} We use this baseline for sparse linear functions (Section \ref{sec:beyond_linear}). This corresponds to $M(P) = \hat{w}^T x_{\text{query}}$, where $\hat{w}$ minimizes the $\ell_1$-norm regularized least squares objective:
\[
\min_{\hat{w}} ~ \frac{1}{2k}|| y - X\hat{w}||_2^2 ~ + ~ \alpha ||\hat{w}||_1.
\]
 We try different values of $\alpha \in \{ 1, 10^{-1}, 10^{-2}, 10^{-3}, 10^{-4} \}$, and report the best solution (achieving the smallest error with $10$ in-context examples) corresponding to $\alpha = 10^{-2}$. To solve the optimization problem, we use the Lasso implementation from Scikit-learn \citep{scikit-learn}
 ~\footnote{\scriptsize\url{https://scikit-learn.org/stable/modules/generated/sklearn.linear_model.Lasso.html}}.

\paragraph{Greedy Tree Learning.} We use this baseline for the class of decision trees. This corresponds to greedily learning a decision tree using the in-context examples, and using it to classify the query input. To construct the tree, at each node (starting from a root node), we choose a coordinate for partitioning the examples into two sets, so as to minimize the variance of $y_i$s in each set, averaged across the two sets. The value associated with a leaf node is the average $y_i$ value of the examples belonging to it. We use Scikit-learn's decision tree regressor \citep{scikit-learn}
~\footnote{\scriptsize\url{https://scikit-learn.org/stable/modules/tree.html\#regression}} implementation for this, with all the arguments set to their default value except the max\_depth argument which is set to 2. 
We considered values $\{1, 2, 3, 4, 5, 6, \text{unbounded}\}$ for the maximum depth and chose the value that performs best at 100 in-context examples which was 2 (which differs from the decision trees sampled from the function class which have depth 4).
We also considered a baseline where we learn this tree using only the signs of each $x_i$ coordinate---after all, the decision tree we are trying to learn depends only on the signs of $x_i$.
In this case, we found the optimal depth to be 4.
 
\paragraph{Tree boosting.} For the class of decision trees, we also consider a tree boosting baseline that corresponds to learning an ensemble of decision trees (see \citet{friedman2001greedy} for a description of the general framework).
Specifically, we use the XGBoost library~\citep{chen2016xgboost}~\footnote{\scriptsize\url{https://github.com/dmlc/xgboost}}, an implementation commonly used for a wide range of real-world machine learning tasks.

We performed a hyperpameter search by considering \{1, 2, 5, 10, 50, 100, 200, 400\} estimators in the ensemble (equivalent to number of boosting rounds), a learning rate of \{0.001, 0.01, 0.1, 0.3, 0.6, 1, 3\}, and a maximum depth of \{1, 2, 3, 4, 6, 10, 16\}.
In general, we found the performance of the learning algorithm to be quite robust.
We chose the hyperparameters obtaining the best performance with 100 training examples, corresponding to 50 estimators, a maximum depth of 4, and a learning rate f 0.1.
We found these hyperparameters to also be optimal when learning based on the signs of each $x_i$.

\paragraph{Learning neural networks with gradient descent.} We use this baseline for the class of two-layer neural networks (Section \ref{sec:beyond_linear}). This corresponds to training a two-layer neural network on the in-context examples, and outputting its prediction on the query point. That is, $M(P) = \hat{f}(x_\text{query})$, where 
\[
\hat{f}(x_\text{query}) = \sum_{i=1}^r ~ \hat{\alpha}_i ~ \sigma(\hat{w}_i^\top x_\text{query}).
\]
Here, $\sigma(\cdot)$ is the ReLU activation.
We find parameters $\hat{\alpha}_i, \hat{w}_i$ by minimizing the squared error of the prediction for the in-context examples
\[
\sum_{i=1}^k \left(\hat{f}(x_i) - f(x_i)\right)^2,
\]
using the Adam optimizer. %\citep{kingma2014adam}.
We use a batch size of 10 (we use full batch when the number of in-context examples is less than 10) with 5000 optimization steps, and set $r = 100$. We use a learning rate of $5\cdot 10^{-3}$ in the case when the data is generated using a neural network, and a learning rate of $5\cdot 10^{-2}$ when the data is generated using a linear function. We consider the setting with $100$ in-context examples, and do a hyperparameter grid search over learning rate $\in \{5\cdot 10^{-4}, 5\cdot 10^{-3}, 5\cdot 10^{-2}, 5\cdot 10^{-1}, 5\}$, $r \in \{100, 400\}$, batch size $\in \{10, 100\}$, optimization algorithm $\in \{\text{adam}, \text{sgd}\}$. All the hyperparameter settings in this grid led to a similar or worse performance compared to the hyperparameter setting we choose. 

%Update the main text hyperparameter discussion also

% over the whole batch of size $k$. We use a learning rate of $5\cdot 10^{-3}$ with 100  optimization steps, and set $r = 100$.
% We performed initial experiments with different values of these parameters (also considering SGD as an optimizer) and did not find a significant improvement.



